\documentclass[11pt,a4paper]{article}
\usepackage[T1]{fontenc}

\usepackage{amsmath}
\usepackage{amssymb}
\usepackage{isabelle,isabellesym}

\usepackage{pdfsetup}

\urlstyle{rm}
\isabellestyle{literal}

% prefer a slightly looser line over text running into the margin
\setlength{\emergencystretch}{3em}


\begin{document}

\title{Work Bounds for Strongly Joinable Balanced Binary Search Trees}
\author{Marco Haucke}
\maketitle

\begin{abstract}
  This entry formalises the complexity analysis of set operations on strongly
  joinable trees as defined by Blelloch et al.~\cite{blelloch2022joinable}.
  It is proved that union, intersection, and difference run in time
  $O(m \log(n/m + 1))$ for input sets of size $m$ and $n$ where $m \le n$.
\end{abstract}

\tableofcontents

\section{Introduction}

Binary search trees are one of the most fundamental data structures in
computer science. Together with their traversals, they can be used to
efficiently implement algorithms for managing ordered data. Balanced
variants such as AVL trees~\cite{adelsonvelskii1962algorithm}, red-black
trees~\cite{guibas1978dichromatic} or weight-balanced
trees~\cite{nievergelt1973bounded} keep the
height logarithmic in the size by maintaining a structural invariant, which
enables simple and efficient implementations of pointwise operations.
Operations that combine two trees are a different matter. Given two sets of
sizes $m \le n$, inserting the elements of the smaller set one by one takes
time $\Theta(m \log n)$, and flattening both trees to lists, merging them and
rebuilding a tree takes $O(m + n)$. Neither is satisfactory across the whole
range of $m$ and $n$.

The \emph{join-based} approach, proposed by Adams~\cite{adams1993efficient},
improves on this.
Here, union, intersection and difference are all expressed through a
single function \texttt{join l k r} that concatenates (i.e. ``joins'') two trees 
$l$ and $r$ and a key $k$ into one balanced tree.

Blelloch, Ferizovic and Sun~\cite{blelloch2022joinable} put Adams' ideas on a
general footing. They characterise a balancing scheme by a real-valued
\emph{rank} function together with a small set of rules relating rank, tree
size, the shape and cost of \texttt{join}, and call
schemes obeying these rules \emph{strongly joinable}. For every strongly
joinable scheme they show that union, intersection and difference of trees
with sizes $m \le n$ take $O\bigl(m \log(\frac{n}{m} + 1)\bigr)$ work, which
is optimal in the comparison model~\cite{hwang1972merging}. 
Further, they prove that AVL trees, red-black trees and weight-balanced trees 
are strongly joinable. 

This topic is not merely of theoretical interest. Adams' original 
code lives on in the SML/NJ Library~\cite{smlnjlib} and in MIT/GNU 
Scheme~\cite{mitscheme}, and the same join-based algorithms underlie the 
\texttt{Set} modules of OCaml's standard library~\cite{ocamlset} as well as 
Jane Street's \emph{Base}~\cite{base}, the immutable \texttt{TreeSet} of 
Scala~\cite{scalatree}, and the \texttt{Data.Set}
implementation of Haskell's \emph{containers} library~\cite{containers}.
Sun et al.'s PAM library~\cite{pamlib,sun2018pam} shows the bound
paying off in practice.

On the Isabelle side, the join-based approach is the subject of
\emph{Functional Data Structures and
Algorithms}~\cite[Chapter~10]{fdsa} and contained inside the \texttt{Set2\_Join} 
locale in \texttt{HOL-Data\_Structures}. Given a \texttt{join} that
satisfies its functional specification, union, intersection and
difference are defined generically and proven correct.

Verified counterparts exist in other proof assistants. The Rocq standard
library's \texttt{MSetAVL}~\cite{msetavl} is a port of OCaml's \texttt{Set}
module, HOL4's \texttt{balanced\_map}~\cite{hol4bmap}, which underlies the
CakeML basis library, is a port of \emph{containers}' \texttt{Data.Map}, and
Breitner et al.~\cite{breitner2018ready} verify \emph{containers}'
\texttt{Data.Set} itself in Coq. All three establish functional correctness
only. On the cost side, Li, Grodin and
Harper~\cite{li2023calf} verify a precise bound for \texttt{join} on
red-black trees in the \emph{calf} framework, but leave the cost of the set
operations to future work. 

To the author's knowledge, this entry contains the first
mechanised proof of the bound. It extends \texttt{Set2\_Join} with a
locale \texttt{StronglyJoinable} that adds the rank-based balance predicate
and the rules of Blelloch et al., and proves within it that the running-time
functions of union, intersection and difference are in
$O\bigl(m \log(\frac{n}{m} + 1)\bigr)$ for input trees of sizes $m \le n$.

Running times are measured by
step-counting functions~\cite[Section~1.5]{fdsa} generated by the
\texttt{time\_fun} command.
As proof of concept, the locale is instantiated
with the AVL trees of \texttt{HOL-Data\_Structures}, as well as red-black trees 
mostly as defined in \texttt{Set2\_Join\_RBT}.

\input{Analytical}

\input{BalancedBinary}

\input{StronglyJoinable}

\input{StronglyJoinableAVL}

\input{StronglyJoinableRBT}

\appendix

\input{Submodularity}

\bibliographystyle{abbrv}
\bibliography{root}

\end{document}
